\documentclass[11pt]{article}
\usepackage[margin=1in]{geometry}
\usepackage{booktabs}
\usepackage{hyperref}

\title{Evaluating Time-Series Foundation Models}
\author{}
\date{}

\begin{document}
\maketitle

\section{Question}

This project evaluates how foundation-model forecasts on the official TIME
test tasks change with target representation, usable covariates, input
normalization, and available context length. Each comparison holds the TIME
dataset windows and metric implementation fixed unless its protocol explicitly
declares another factor.

\section{Current migrated comparisons}

The executable surface contains a foundation benchmark of Chronos-2,
Chronos-Bolt, TS-ICL, and Seasonal Naive; TimesFM-3 is excluded from the active
experiment set. It also contains a Chronos-2 comparison between native
multivariate targets, independent univariate targets, and other target
histories supplied as past-only covariates. Seasonal Naive is deterministic
and provides the taskwise MASE
normalizer. It also fixes the evaluation grid to cells with finite target
support, finite Seasonal Naive predictions on that support, and finite
Seasonal Naive MASE. Every learned model uses this same grid; a non-finite
forecast on expected support invalidates its task rather than changing the
aggregation support. Inference timing excludes loading, dataset construction,
metric computation, and artifact saving.

Two further executable comparisons use the same three learned backbones and
evaluation grid. The maximum-context study evaluates five values per backbone,
halving the model-specific maximum four times. Its summary aggregates
taskwise MASE divided by the matched Seasonal Naive MASE for every model,
context-size, and horizon-size cell. The normalization study compares the
unchanged input with a per-window, per-variate z-score computed after context
truncation. Constant inputs use unit scale, and every output quantile is
transformed back to the target units before evaluation.

\section{Evidence boundary}

Current result claims are limited to project-owned manifests and aggregates
produced under the shared Seasonal evaluation grid. Covariate generalization,
normalization, and context-size studies make no empirical claim until their
new launch artifacts are synchronized and analyzed.

\end{document}
